\documentclass{article}
\usepackage{graphicx} % Required for inserting images

% Added stuff
\usepackage{amsmath}
\usepackage{amssymb}
\newcommand{\N}{\mathbb{N}}
\newcommand{\Z}{\mathbb{Z}}
\newcommand{\Q}{\mathbb{Q}}
\newcommand{\R}{\mathbb{R}}
\newcommand{\C}{\mathbb{C}}
\newcommand{\bigp}[1]{\left(#1\right)}
\newcommand{\bigb}[1]{\left[#1\right]}
\newcommand{\bigabs}[1]{\left|#1\right|}
\newcommand{\summation}[1][0]{\sum_{n=#1}^\infty}
\newcommand{\summationk}[1][0]{\sum_{k=#1}^\infty}
\newcommand{\summationm}[1][0]{\sum_{m=#1}^\infty}

\title{Infinite Series 9}
\author{Siddharth K}
\date{March 2026}

\begin{document}

\maketitle

\section*{Problem 5}
Let
$$x = \sqrt[p]{k + \sqrt[p]{k + \sqrt[p]{k + \cdots}}}.$$
Then
$$x = \sqrt[p]{k + x},$$
$$x^p - x - k = 0.$$
So $x$ is the root of $x^p - x - k = 0$.

Let $x_1 = \sqrt[p]{k}$, $x_{n + 1} = \sqrt[p]{k + x_n}$. Then the sequence $x_n$ converges because it's increasing and bounded. It's increasing because
$x_1 = \sqrt[p]{k} < x$, and for any $x_n < x$,
$$x^p - x - k < 0 \implies x^p < k + x \implies x < \sqrt[p]{k + x},$$
so inductively, $x_n < x_{n + 1} = \sqrt[p]{k + x}$.

The sequence is bounded because the function $f(x) = \sqrt[p]{k + x}$ is increasing and if $x_n \le x$ then
$$x_{n + 1} = f(x_n) \le f(x) = x,$$
and $x_1 < x$ so inductively, $x_n \le x$ for all $n$.

\section*{Problem 6}
Let
$$x = \sqrt[n]{q + p\sqrt[n]{q + p\sqrt[n]{q + \cdots}}}.$$
Then
$$x = \sqrt[n]{q + px},$$
$$x^n - px = q,$$
as desired. Now, let
$$y = \sqrt[m/n]{q - p\sqrt[m/n]{q - p\sqrt[m/n]{q - \cdots}}}.$$
Then
$$y^n = \sqrt[m]{q - py}, \quad y^{nm} - py = q.$$
The desired nested radical is equal to $y^n$, so it satisfies
$$x^m - px^{1/n} = q.$$

\section*{Problem 9}
Let
$$f(x) = \sqrt{x + \sqrt{x + 1 + \cdots}}.$$
Then
$$f(x) = \sqrt{x + f(x + 1)},$$
$$f(x)^2 = x + f(x + 1).$$
So the nested radical is equal to $f(1)$ for some $f$ that satisfies the functional equation.

\section*{Problem 10}
Let
$$f(x) = \sqrt[3]{6x - 2 + (2x - 1)^2 \sqrt[3]{(6x + 4) + \cdots}}.$$
Then
$$f(x)^3 = 6x - 2 + (2x - 1)^2f(x + 1).$$
If $f$ is a polynomial, it must be linear with leading coefficient $2$, so we guess that $f(x) = 2x$. The functional equation holds, so the radical is $f(1) = 2$.

\section*{Problem 13}
Given a cycle of length $n$, $f^n(x_1) = x_1$, so $x_1$ is a fixed point of $f^n$. So stability should be determined using the derivative of $f^n$. Using the chain rule,
$$(f^n)'(x_1) = f'(x_n)f'(x_{n - 1}) \cdots f'(x_1).$$
So stability of a cycle can be defined as: if $|(f^n)'(x_1)| < 1$, the cycle is attracting. If it's $> 1$, the cycle is repelling. If it's $= 1$, the cycle is neutral.

\section*{Problem 15}
Let $L = x^{x^{\cdots}}$. Then $L = x^L$, or
$$L = e^{L\log x}.$$
Also,
$$Le^{-L\log x} = 1,$$
$$(-L\log x)e^{-L\log x} = -\log x.$$
By definition of the Lambert $W$ function, $-L\log x = W(-\log x)$, so
$$L = -\frac{W(-\log x)}{\log x}.$$

\section*{Problem 16}
The power tower converges whenever
$$e^{-e} \le \frac{1}{x} \le e^{1/e},$$
so
$$e^{-1/e} \le x \le e^e.$$
The value of the power tower is
$$-\bigp{\frac{W(-\log (1/x))}{\log (1/x)}}^{-1} = \frac{\log x}{W(\log x)}.$$

\section*{Problem 17}
For each part, we test whether
$$W(x)\exp(W(x)) = x.$$
\subsection*{Part (a)}
$$\frac{i\pi}{2}e^{i\pi/2} = \frac{i\pi}{2}\cdot i = -\frac{\pi}{2}.$$
\subsection*{Part (b)}
$$-1 \cdot e^{-1} = -\frac{1}{e}.$$
\subsection*{Part (c)}
$$\log 2 \cdot e^{\log 2} = 2\log 2.$$

\end{document}